% ============================================================================
\clearpage
\section{Producto y pregunta de investigación}\label{sec:producto}
% ============================================================================

\ficha{Cadenas de opciones con bid/ask, precio del subyacente, posiciones actuales, calendario de
resultados y dividendos.}
{Medir $\Q$, describir su cambio, pronosticar $\P$ y comparar mantener con modificar.}
{Recomendación con bandas de incertidumbre, o abstención motivada; registro reproducible.}
{Separación estricta entre $\Q$, $\P$ y acción; ninguna decisión sin calibración ni riesgo conjunto.}

\subsection{Qué hace el sistema en cada sesión}

El producto responde, al inicio de cada sesión, una pregunta operativa concreta: \emph{dado lo que
descuentan hoy las opciones y lo que ya se sabe por el precio, ¿conviene mantener la posición o
modificarla?} Para llegar a esa respuesta encadena cinco pasos, cada uno con su propia salida
verificable:

\begin{enumerate}
\item \textbf{Reconstruir lo que descuentan las opciones.} A partir de un \emph{snapshot} de
cotizaciones se obtiene una superficie de volatilidad sin arbitraje estático y, de ella, la
distribución implícita $\Q_t$ del rendimiento a plazo constante, con bandas que reflejan el ancho de
los spreads.
\item \textbf{Cuantificar cómo cambió.} Se compara $\Q_t$ con $\Q_{t-1}$ mediante distancias entre
funciones cuantil, cambios firmados por percentil y factores funcionales filtrados.
\item \textbf{Estimar si esa información mejora el pronóstico.} Los cambios --- en especial su parte
no explicada por la información base --- entran como variables candidatas a modelos que pronostican la
distribución física $\P$ de rendimientos futuros, evaluados fuera de muestra.
\item \textbf{Comparar mantener con modificar.} Con escenarios conjuntos de la cartera se evalúa el
rendimiento esperado, el CVaR y los costos de cada alternativa; los cambios pequeños que no superan
costos e incertidumbre se descartan.
\item \textbf{Registrar.} Cada resultado conserva las entradas, versiones, parámetros, hora de corte y
motivo de abstención, de modo que la corrida pueda repetirse exactamente.
\end{enumerate}

\subsection{Hipótesis principal y cómo se contrasta}\label{sec:hipotesis}

Sea $Y_{t,h}$ una variable objetivo medida entre el instante de decisión $t$ y $t+h$ (rendimiento,
pérdida o volatilidad realizada), $\mathcal{I}^{base}_t$ la información ya observable en precio,
volatilidad y mercado, y $\epsilon_t$ la innovación de la distribución implícita definida en la
\secref{sec:innovacion}. Con una regla de puntuación propia $S$ (menor es mejor) y dos pronósticos de
la \emph{misma familia} de modelos, uno con $\mathcal{I}^{base}_t$ y otro con
$\mathcal{I}^{base}_t\cup\{\epsilon_t\}$, se define el diferencial de pérdida
\begin{equation}
d_{t,h}=S\bigl(\widehat F^{\,base}_{t,h},\,Y_{t,h}\bigr)-S\bigl(\widehat F^{\,base+\epsilon}_{t,h},\,Y_{t,h}\bigr).
\label{eq:diferencial}
\end{equation}
La hipótesis de trabajo es $H_1:\ \E[d_{t,h}]>0$ fuera de muestra, frente a $H_0:\ \E[d_{t,h}]\le 0$.
Se evalúa por separado para cada objetivo y horizonte, con intervalos que respetan la dependencia
temporal (\secref{sec:evidencia}).

\begin{noconcluir}
\begin{itemize}
\item Una mejora en el pronóstico del \emph{riesgo} no implica una mejora en el \emph{signo} del
rendimiento, ni al revés; ambas conclusiones son útiles y se reportan por separado.
\item Una mejora en el pronóstico no implica una mejor política de posición: puede no pagar sus
costos. Son dos pruebas distintas (\secref{sec:dos_pruebas}).
\item La hipótesis no presupone una dirección rentable; esa dirección, si existe, se mide.
\end{itemize}
\end{noconcluir}

\subsection{Alcance de la primera versión}

La primera versión se limita a SPY y entre tres y cinco acciones líquidas, con posiciones en
acciones o ETF y sin apalancamiento como supuesto inicial de investigación. Gestionar opciones como
posición exigiría una fase específica de valoración y escenarios conjuntos (\secref{sec:opciones}).

\begin{table}[H]
\centering\small
\caption{Alcance de la primera versión.}
\label{tab:alcance}
\begin{tabularx}{\linewidth}{@{}l l Y@{}}
\toprule
\encab{Elemento} & \encab{Estado} & \encab{Comentario}\\
\midrule
Universo & Propuesto & SPY y 3--5 acciones líquidas; lista definitiva pendiente.\\
Instrumentos gestionados & Supuesto inicial & Acciones y ETF; opciones como fase aparte.\\
Apalancamiento & Supuesto de investigación & Sin apalancamiento en la primera versión.\\
Capital de referencia & Pendiente & 50,000; falta definir la moneda.\\
Posiciones actuales & Pendiente & Necesarias para evaluar ``mantener frente a modificar''.\\
Horizonte de decisión & Pendiente & Candidatos: 1, 5 y 20 sesiones.\\
Presupuesto de datos & Pendiente & Condiciona la fuente y la historia disponible.\\
\bottomrule
\end{tabularx}
\end{table}

\begin{pendiente}
Universo definitivo, moneda de los 50,000, posiciones, horizonte y presupuesto de datos. No son
necesarios para \emph{diseñar} el motor, pero sí para \emph{parametrizar} su uso: sin ellos no se
pueden fijar límites de concentración, costos ni el benchmark comparable.
\end{pendiente}

\subsection{Tres horizontes que no deben mezclarse}\label{sec:horizontes}

El sistema maneja tres escalas de tiempo distintas (\figref{fig:horizontes}). La \textbf{frecuencia
de observación} es cada sesión, incluidas las ventanas de apertura. El \textbf{horizonte de la
distribución implícita} es el plazo de las opciones que se interpolan; inicialmente 30 días
calendario y constante, de modo que en $t+1$ se sigue midiendo a 30 días y no a 29. El
\textbf{horizonte de decisión y pronóstico} es aquel sobre el que se evalúa el rendimiento de una
decisión, por ejemplo 1, 5 y 20 sesiones.

\begin{figure}[H]
\centering
\resizebox{\linewidth}{!}{% Tres horizontes: observación, distribución implícita y decisión/pronóstico.
\begin{tikzpicture}[x=0.36cm,y=1cm,nota/.style={font=\sffamily\scriptsize,text=tinta2,align=left}]
% Etiquetas de fila
\node[titulo,anchor=east,align=right] at (-0.8,0) {Frecuencia de\\ observación};
\node[titulo,anchor=east,align=right] at (-0.8,-1.35) {Horizonte de la\\ distribución $\Q$};
\node[titulo,anchor=east,align=right] at (-0.8,-2.95) {Horizonte de decisión\\ y pronóstico $\P$};
% Fines de semana sombreados
\foreach \s in {5,12,19,26} \fill[grisclaro] (\s-0.5,0.45) rectangle (\s+1.5,-3.75);
% Instante de decisión
\draw[tinta,line width=0.6pt] (0,0.55) -- (0,-3.75);
\node[nota,anchor=south] at (0,0.55) {instante de decisión $t$};
% Observación: cada sesión hábil
\foreach \d in {0,...,30}{%
  \pgfmathtruncatemacro{\w}{mod(\d,7)}%
  \ifnum\w<5 \fill[tenue] (\d,0) circle (2.1pt);\fi}
\node[nota,anchor=west] at (31.2,0) {cada sesión: apertura\\ y cierre, sin operar\\ necesariamente};
% Distribución implícita a plazo constante
\draw[line width=6pt,qazul!55] (0,-1.15) -- (30,-1.15);
\draw[line width=6pt,qazul!22] (1,-1.55) -- (31,-1.55);
\node[nota,anchor=west] at (31.2,-1.15) {$\tau=30$ días calendario};
\node[nota,anchor=west] at (31.2,-1.62) {en $t+1$ el plazo sigue\\ siendo 30 días (interpolado)};
% Horizontes de decisión
\draw[flecha,draw=qnaranja,line width=0.9pt] (0,-2.45) -- (1,-2.45);
\node[nota,anchor=west] at (1.3,-2.45) {1 sesión};
\draw[flecha,draw=qnaranja,line width=0.9pt] (0,-2.95) -- (7,-2.95);
\node[nota,anchor=west] at (7.3,-2.95) {5 sesiones};
\draw[flecha,draw=qnaranja,line width=0.9pt] (0,-3.45) -- (28,-3.45);
\node[nota,anchor=west] at (28.3,-3.45) {20 sesiones};
% Eje de tiempo
\draw[guia] (-0.5,-3.9) -- (31,-3.9);
\foreach \d in {0,7,14,21,28}{\draw[guia] (\d,-3.9) -- (\d,-4.0);
  \node[nota,anchor=north] at (\d,-4.02) {\d};}
\node[nota,anchor=north] at (15,-4.35) {días calendario desde $t$ (sombreado: fines de semana)};
\end{tikzpicture}
}
\caption{Los tres horizontes. Observar cada sesión no obliga a operar cada sesión; la distribución
implícita se mide siempre a 30 días calendario; los pronósticos se evalúan a 1, 5 y 20 sesiones.}
\label{fig:horizontes}
\end{figure}

La distinción importa porque una distribución a 30 días \emph{no se convierte por simple escalado} en
una distribución a una sesión. La varianza no se acumula de manera uniforme: los fines de semana
aportan poco, un día de resultados trimestrales aporta mucho, y la estructura temporal de
volatilidades rara vez es plana. Las figuras~\ref{fig:acumulacion} y~\ref{fig:riesgo_sesion}
lo ilustran con un calendario sintético de 30 días que incluye un anuncio de resultados.

\grafica{s1_acumulacion_varianza}{Acumulación de varianza en 30 días calendario (datos sintéticos).
La escalera muestra una acumulación plausible: casi nada en fines de semana y un salto el día de
resultados. La recta es lo que supone un escalado lineal en días calendario.}{fig:acumulacion}

\grafica{s1_riesgo_sesion}{Desviación estándar de una sesión (datos sintéticos). Las dos primeras
barras son los valores del mundo sintético para una sesión normal y para la sesión de resultados; las
dos últimas, lo que se obtiene escalando la varianza a 30 días por días calendario o por sesiones.
Ningún escalado recupera el riesgo de la sesión de resultados.}{fig:riesgo_sesion}

\begin{ideaclave}
El plazo de $\Q$ describe qué se mide; el horizonte de decisión describe qué se evalúa. Pasar de uno a
otro exige un modelo explícito (calendario de eventos, estructura temporal, supuestos sobre fines de
semana), nunca una regla de raíz del tiempo aplicada por defecto.
\end{ideaclave}
